\scanpage{54}

\uline{Lemme}
\diagram{\begin{tikzcd}
\mathcal M\arrow[r,shift left,"\lambda_!"]&\deleted{\unsure{A}}\mathrm{Ch}_\bullet(A)\arrow[l,shift left,"\lambda^*"]
\end{tikzcd}}

\deleted{$\Lambda$ foncteur \uline{exact}, \inserted{\unsure{Supposé}} transformant}

\deleted{de ce} $\mathcal M,A$ cat. abéliennes

$\lambda_!$ foncteur exact, transformant projectifs en projectifs

\deleted{i.e. pour $P$ projectif\par
\[
\begin{array}{ll}
a)&\lambda_{!i}(P)\text{ projectif}\ \deleted{\text{pour }P\text{ projectif}}\quad i\in\mathbb N\\
b)&H_i(\lambda_!(P))=0\quad\text{si }i>0\\
c)&H_0(\lambda_!(P))\text{ projectif}
\end{array}
\]}

$\lambda^*$ \deleted{foncteur} $(\lambda_!,\lambda^*)$ couple de foncteurs adjoints,
\deleted{$\lambda^*$ exact \uline{à droite}}
\inserted{\[
\left\{\begin{array}{ll}
a)&\lambda_!\text{ transforme mono en mono (donc exact)}\\
b)&\lambda^*\text{ transforme épi en épi (donc exact)}
\end{array}\right.
\]}

\[
[\Longrightarrow\lambda_!\text{ transforme projectifs en projectifs}]
\]
\[
\varphi=H_0\circ\lambda_!
\]

\deleted{\unsure{$A$}}
\margin{\uuline{\unsure{Réécrivons}}}

\uline{Lemme 1}
\diagram{\begin{tikzcd}
\mathcal M\arrow[r,"{\deleted{\unsure{\varphi_\bullet}}\Lambda_\bullet}"]\arrow[dr,dashed,"L\varphi"']&\mathrm{Ch}_\bullet(A)\arrow[d,"c=\mathrm{can}."]\\
&\mathrm D_\bullet(A)
\arrow[ul,phantom,"\simeq" description]
\end{tikzcd}}

foncteur \uline{exact}, transformant \uline{projectifs en \unsure{idem}}
\margin{(plus gén, $H_i(\Lambda_\bullet(P))=0$ si $P$ proj, $i>0$}

\[
\varphi_0=\underline{H_0}\circ\Lambda=(x\mathbin{-}H_0(\Lambda_\bullet(x)))
\]

Supposons qu'il y ait assez de projectifs dans $\mathcal M$.

Alors pour \deleted{tout} $x\in\mathcal M$,
\[
(L\varphi)(x)\simeq c\Lambda_\bullet(x)
\]

\uline{Dém.} \deleted{\unsure{$P_\bullet$}}
$P_\bullet^x\longrightarrow x$ résolution projective,
\[
(L\varphi)(x)\simeq\varphi_0(P_\bullet^x)
\]
\deleted{\unsure{$\forall i$}} Considérons le bicomplexe
\diagram{\begin{tikzcd}
\varphi_0(P_\bullet^x)&\\
\Lambda_\bullet(P_\bullet^x)\arrow[u,"\text{qu. iso}"]\arrow[r]&\Lambda_\bullet(x)
\end{tikzcd}}
\margin{car $\forall i$, $\Lambda_\bullet(P_i^x)$ \uline{résolution} de $\varphi(P_i^x)$}
qu. iso, car pour tt $j$, $\Lambda_j(P_\bullet^x)$
\deleted{\unsure{$\longrightarrow$}} résolution de $\Lambda_j(x)$
($\Lambda_j$ \uline{exact})

\margin{\figure{\begin{tikzpicture}
\node[draw,ellipse,rotate=50,inner sep=7pt] {\unsure{Utiliser} \unsure{Notations} $\lambda_!,\lambda^*$};
\end{tikzpicture}}}

\uline{Lemme 2} Supposons que $\Lambda_\bullet=\lambda_!$
\deleted{soit exact et}\inserted{ait un} adjoint à droite
\deleted{exact} $\lambda^*$ et $\lambda_!$ transforme mono en mono,
$\lambda^*$ épi en épi (donc $\lambda_!,\lambda^*$ exacts).
Alors $\lambda_!$ transforme projectifs en projectifs,
$c\circ\lambda_!\simeq L\varphi$ (où $\varphi=H_0\circ\lambda_!$),
\deleted{et pour que $\lambda_!$}
Posons
\[
W_{\mathcal M}=\Lambda_\bullet^{-1}(W_{\mathrm{Ch}_\bullet(A)})\ ;
\]
\deleted{Alors \unsure{conditions équivalentes}}

\mspage{53}
